\section{Conclusion}
\label{sec:conclusion}

We developed a representation of the \acs{drausp} as a finite-horizon \acs{mdp}, which enables us
to use reinforcement learning techniques to learn acceptance and scheduling policies. For optimization, we used
a value-based deep \acs{rl} approach, specifically Q-learning with a deep neural network as function approximator.
We constructed a stable baseline approach in Section \ref{sec:rl-approach} using a Standard \acs{dqn}. We
tested the behaviour of the learned $Q$-function with respect to monotonicity relations inherent to the problem.
Doing so, we constructed a monotonicity test to evaluate the learned $Q$-function (see \ref{subsec:monotonicity-test}).\\
As a second approach, we took advantage of the known monotonicity relations presented in \ref{subsec:monotonicity} 
and constructed a monotonicity-aware \acs{dqn}, using lattice layers. We built up a deep lattice
network architecture that enforces monotonicity by design, and again evaluated its performance 
and monotonicity behaviour on the same benchmark instances as the Standard \acs{dqn}.\\
We saw that even without any knowledge about the problem structure at the beginning, the
Standard \acs{dqn} still manages to learn monotonicity. The Lattice-\acs{dqn} 
on the other hand, enforces monotonicity by design and therefore satisfies all monotonicity tests exactly. 
But using the Lattice-\acs{dqn} comes with a trade-off: the more constrained function class of the Lattice-\acs{dqn}
makes it harder to learn the optimal policy. This leads to a slower convergence and a slightly worse performance in the end.\\

Future work could focus on improving the performance of the Lattice-\acs{dqn} by using a more complex architecture, e.g. by
adding more layers or increasing the number of lattices per layer. Another approach could be to enforce monotonicity only on a subset of the features.
In this case one could use a hybrid architecture,
where non-monotonic features are first processed by standard feedforward layers, while the monotonicity-relevant features
bypass these layers (or only pass through monotone transformations) and are combined with the resulting representation only
at the final lattice layers. This would allow the network to learn a more complex function from the non-monotonic features
while still enforcing end-to-end monotonicity with respect to the relevant inputs. Note that simply stacking feedforward
layers before lattice layers on all inputs would not preserve this guarantee, since composing a monotone function with a
non-monotone one is generally not monotone.
In our numerical experiments the Standard \acs{dqn} still manages to learn monotonicity, but it is not guaranteed to do so.
One could test if this holds true if the network architecture is changed or instances get larger and more complex.